\documentclass[11pt,a4paper]{article}
\usepackage[utf8]{inputenc}
\usepackage[T1]{fontenc}
\usepackage{amsfonts}
\usepackage[french]{babel}
\frenchbsetup{StandardLists=true}
\usepackage{enumitem}
\usepackage{amssymb}
\usepackage{amsmath}
\usepackage[left=2cm,right=2cm,top=2cm,bottom=2cm]{geometry}
\usepackage{multirow}
\usepackage[dvipsnames]{xcolor}
\usepackage{fouriernc}
\usepackage{graphicx}
\usepackage{setspace}
\singlespacing
\setlength{\parindent}{0pt}
\usepackage{multicol}
\setlength{\columnseprule}{1pt}
\usepackage{tikz-osci}
\usepackage{fancyhdr}
\pagestyle{fancy}
\renewcommand\headrulewidth{1pt}
\fancyhead[L]{Lycée Denis Diderot }
\fancyhead[R]{Classe de Première}
\setcounter{page}{2}\fancyfoot[C]{page \thepage}
\fancyfoot[R]{Année 2025-2026}
\fancyfoot[L]{Alexis RUIZ - CORRIGÉ}
\usepackage{multirow,tabularx,booktabs}
\newcommand{\cadreblanc}[1]{%
\medskip\framebox[\linewidth][c]{%
\rule{0mm}{#1mm}}\par
\medskip}
\renewcommand{\arraystretch}{2}
\newcommand*\acocher{\quitvmode{\fboxsep0pt \fboxrule0.8pt \fbox{\vrule height1.8ex depth.1ex width0pt \vrule height0pt depth0pt width1.9ex }}}
\usepackage{array}
\AtBeginDocument{\shorthandoff{;}}
\begin{document}
\begin{center}\LARGE{Contrôle 1 - Lecture d'oscillogrammes - CORRIGÉ}\end{center}\vspace{0.3cm}
\normalsize
\hrule\vspace{0.2cm}
\flushleft \textbf{NOM :\\[0.2cm]
Prénom:\\[0,2cm]}
\hrule\vspace{0.2cm}
\vfill 
\textbf{Consigne :} Lire sur l'oscillogramme les valeurs de U\textsubscript{max} et T, puis calculer les valeurs de U (Arrondir à 0,1) et \textit{f} (Arrondir à l'unité)
\vfill

% Oscillogramme 1 - 15V, T=0.08s, 5V/div, 20ms/div
\begin{minipage}{0.45\textwidth}
\centering
\osci[%
scale=0.75,
second channel=0,
screen offset one=0,
time div=20,
voltage div one=5,
sample rate=200,
xy mode=0,
func one=15*sin(2*180/0.08*x),
color one=003399,
graph back color=FFFFFF,
info back color=D6D6D6,
info text color=000000,
main axis color=000000,
grid color=CCCCCC
]

\vspace{0.1cm}

\begin{tabular}{|m{3.2cm}|m{3.4cm}|}
\hline
$\mathrm{U}_{\text{max}} = \textcolor{red}{15\text{ V}}$ & $\mathrm{T} = \textcolor{red}{80\text{ ms}=0,08 \text{ s}}$ \\
\hline
$\mathrm{U} = \textcolor{red}{10,6\text{ V}}$ & $\mathrm{f} = \textcolor{red}{13\text{ Hz}}$ \\
\hline
\end{tabular}
\end{minipage}
\hfill
% Oscillogramme 2 - 32V, T=0.02s, 10V/div, 10ms/div
\begin{minipage}{0.45\textwidth}
\centering
\osci[%
scale=0.75,
second channel=0,
screen offset one=0,
time div=10,
voltage div one=10,
sample rate=200,
xy mode=0,
func one=32*sin(2*180/0.02*x),
color one=CC3300,
graph back color=FFFFFF,
info back color=D6D6D6,
info text color=000000,
main axis color=000000,
grid color=CCCCCC
]

\vspace{0.1cm}

\begin{tabular}{|m{3.2cm}|m{3.4cm}|}
\hline
$\mathrm{U}_{\text{max}} = \textcolor{red}{32\text{ V}}$ & $\mathrm{T} = \textcolor{red}{20\text{ ms}=0,02 \text{ s}}$ \\
\hline
$\mathrm{U} = \textcolor{red}{22,6\text{ V}}$ & $\mathrm{f} = \textcolor{red}{50\text{ Hz}}$ \\
\hline
\end{tabular}
\end{minipage}

\vspace{1cm}

% Oscillogramme 3 - 19V, T=0.025s, 5V/div, 5ms/div
\begin{minipage}{0.45\textwidth}
\centering
\osci[%
scale=0.75,
second channel=0,
screen offset one=0,
time div=5,
voltage div one=5,
sample rate=200,
xy mode=0,
func one=19*sin(2*180/0.025*x),
color one=009900,
graph back color=FFFFFF,
info back color=D6D6D6,
info text color=000000,
main axis color=000000,
grid color=CCCCCC
]

\vspace{0.1cm}

\begin{tabular}{|m{3.2cm}|m{3.4cm}|}
\hline
$\mathrm{U}_{\text{max}} = \textcolor{red}{19\text{ V}}$ & $\mathrm{T} = \textcolor{red}{25\text{ ms}=0,025 \text{ s}}$ \\
\hline
$\mathrm{U} = \textcolor{red}{13,4\text{ V}}$ & $\mathrm{f} = \textcolor{red}{40\text{ Hz}}$ \\
\hline
\end{tabular}
\end{minipage}
\hfill
% Oscillogramme 4 - 5.6V, T=0.012s, 2V/div, 5ms/div
\begin{minipage}{0.45\textwidth}
\centering
\osci[%
scale=0.75,
second channel=0,
screen offset one=0,
time div=5,
voltage div one=2,
sample rate=200,
xy mode=0,
func one=5.6*sin(2*180/0.012*x),
color one=990099,
graph back color=FFFFFF,
info back color=D6D6D6,
info text color=000000,
main axis color=000000,
grid color=CCCCCC
]

\vspace{0.1cm}

\begin{tabular}{|m{3.2cm}|m{3.4cm}|}
\hline
$\mathrm{U}_{\text{max}} = \textcolor{red}{5,6\text{ V}}$ & $\mathrm{T} = \textcolor{red}{12\text{ ms}=0,012 \text{ s}}$ \\
\hline
$\mathrm{U} = \textcolor{red}{4,0\text{ V}}$ & $\mathrm{f} = \textcolor{red}{83\text{ Hz}}$ \\
\hline
\end{tabular}
\end{minipage}

\vfill 

\end{document}
